\label{sec:huntington-hill}

\subsection{Priority Formula}

The Huntington-Hill priority for the $(s+1)$-th seat (when a state
currently has $s$ seats) is:
\[
  P_i^{\text{HH}}(s) = \frac{p_i}{\sqrt{s(s+1)}}.
\]
For exact arithmetic, this is represented as:
\[
  P_i^{\text{HH}}(s) = \frac{p_i^2}{s(s+1)},
\]
which eliminates the square root by squaring both numerator and denominator
(preserving the comparison order since both quantities are positive).
The implementation is:

\begin{verbatim}
fn hh_priority(pop: u128, seats: u32) -> Priority {
    Priority {
        numerator: pop * pop,
        denominator: seats as u128 * (seats as u128 + 1),
    }
}
\end{verbatim}

\subsection{Overflow Analysis}

For 2020 Census data, the maximum population is California:
$p_{\text{CA}} = 39{,}538{,}223$.
The maximum numerator is $p_{\text{CA}}^2 = (3.95 \times 10^7)^2
\approx 1.56 \times 10^{15}$.
The maximum denominator is $s \times (s+1)$ where $s$ is at most 52
(California's seat count): $52 \times 53 = 2{,}756$.
The cross-multiplication value in the comparison is at most
$p_i^2 \times s_j(s_j+1) \leq 1.56 \times 10^{15} \times 2756
\approx 4.3 \times 10^{18}$,
well within \texttt{u128}'s capacity of approximately $3.4 \times 10^{38}$.

\subsection{Tie-Breaking}

Ties occur when two states have identical priority values, i.e.,
when $p_i^2 / (s_i(s_i+1)) = p_j^2 / (s_j(s_j+1))$.
In exact integer arithmetic, this requires $p_i^2 \times s_j(s_j+1) =
p_j^2 \times s_i(s_i+1)$, which for census-scale populations is
extremely rare.
When ties occur, the implementation breaks them alphabetically by
state name (a deterministic and documented rule).
The Census Bureau's own computation uses the same Huntington-Hill
priority formula; any tie-breaking rule is valid as long as it
is documented, since true ties in priority values are mathematically
equivalent and either state is equally entitled to the seat.

\subsection{The Last Seat for 2020}

The 385th seat (the last seat allocated above the constitutional floor)
in the 2020 apportionment was awarded to New York.
The priority value at this step was:
\[
  P_{\text{NY}}(25) = \frac{20{,}201{,}249^2}{25 \times 26}
  = \frac{4.08 \times 10^{14}}{650}
  \approx 6.28 \times 10^{11}.
\]
The next highest priority was for a competing state (Montana or
another state near its rounding boundary), with a priority below
New York's.
The Census Bureau publishes the full priority sequence, and
\textsc{bisect-apportion}'s output matches it.
